\section{An extension criterion for stabilizer codes}\label{sec:extension}

\subsection{Pauli space, distance, and degeneracy}

All vector spaces are over \(\F\). Write
\[
 V_n=\F^{2n}=\{(x\mid z):x,z\in\F^n\},\qquad
 \langle(x\mid z),(x'\mid z')\rangle=x\cdot z'+z\cdot x'.
\]
At each coordinate, the pairs \((0,0),(1,0),(0,1),(1,1)\) represent
\(I,X,Z,Y\), respectively, with scalar phases omitted. Vector addition
then represents multiplication of Pauli operators up to phase. The
symplectic product is zero for commuting operators and one for
anticommuting operators. The support of a word is the set of coordinates
where its letter is not \(I\).
The Pauli weight is \(\wt(x\mid z)=|\{i:(x_i,z_i)\ne(0,0)\}|\).
A \(Z\)-only word has letters \(I\) and \(Z\), so its \(X\)-component vanishes.
A stabilizer is represented by its binary subspace; Pauli phases are omitted.
A subspace is isotropic if every pair of its words has symplectic product zero.
A binary \(\qc{n,k}\) stabilizer is an isotropic subspace
\(S\leq V_n\) of dimension \(n-k\). Its centralizer is \(D=S^\perp\),
of dimension \(n+k\). For a Pauli stabilizer this commuting space is
also called its normalizer. For \(k>0\) its distance is
\[
 d(S)=\min\{\wt(w):w\in S^\perp\setminus S\}.
\]
The encoded states are the common eigenspace of independent commuting
generators. The centralizer contains the Pauli operators commuting with
all these generators; its elements outside \(S\) act nontrivially on
the encoded information. Stabilizer words themselves act trivially there,
so a low-weight stabilizer word is permitted. A code is pure to distance
\(d\) when every nonzero word of its centralizer has weight at least \(d\).
For a code of distance at least \(d\), degeneracy at that threshold
means that \(S\) contains a nonzero word of weight below \(d\).
An additive code here is an \(\F\)-linear subspace of Pauli space;
it need not be isotropic. For an additive code \(C\), \(d(C)\) means the
minimum weight of a nonzero codeword, with \(d(\{0\})=+\infty\).
The notation \((\ell,2^c,d)_4\) means an additive code of length
\(\ell\), binary dimension \(c\), and minimum distance \(d\) over
the four-letter Pauli alphabet.

Write \(A_j\) for the number of stabilizer words of weight \(j\); their
list is the weight enumerator. Two additive codes are physically equivalent
if a permutation of coordinates and independent permutations of
\(X,Y,Z\) at each coordinate carry one onto the other. These relabelings
preserve weight and commutation.

Split the coordinates into \(T\), of size \(m\), and its complement,
of size \(f=n-m\). Write \(V_n=V_T\oplus W\), with \(W=V_f\).
We view \(V_T\) and \(W\) as subspaces of \(V_n\) by padding with
identity letters on the other coordinates. Write \(\pi_T\) and
\(\pi_W\) for the two coordinate projections. A prefix is a word on
\(T\), and a suffix is a word on its complement.

\begin{center}\small
\begin{tabularx}{\linewidth}{@{}lX@{}}\toprule
Symbol&Meaning\\\midrule
\(n,k,d\)&Physical qubits, encoded qubits, and target distance\\
\(S,D=S^\perp,A_j\)&Stabilizer, centralizer, and stabilizer weight counts\\
\(T,m,f\)&Selected support, its size, and complementary size \(n-m\)\\
\(R,r\)&Stabilizer supported on \(T\), and its dimension\\
\(C,P=C^\perp\)&Shortened centralizer, and its dual on the complement\\
\(E,a,b\)&Local quotient \(R_T^\perp/R\), its dimension, and \(\dim P\)\\
\(M,\tau,\Gamma\)&Extension map, its symplectic adjoint, and the Gram matrix on \(P\)\\
\(K,q,t\)&Kernel of \(M\), its dimension, and its restricted form rank\\
\(N,h\)&Kernel of \(\tau\), and hull dimension \(\dim(C\cap C^\perp)\)\\
\(H,Q\)&Local centralizer \(S^\perp\cap V_T\), and syndrome image \(\im M\)\\
\(\sigma,c_R,\delta\)&Syndrome map, local coset minimum weight, and complementary coset minimum weight\\
\(\varepsilon,\mathcal{K}_j^\ell\)&Weight-parity function and quaternary Krawtchouk polynomial\\
\(\mu\)&Half the local quotient dimension in the search appendix\\\bottomrule
\end{tabularx}
\end{center}

\subsection{The local quotient and the shortened code}

The local data have different roles. The space \(R=S\cap V_T\)
records stabilizer words supported on \(T\); the shortened centralizer
\(C=S^\perp\cap W\) records centralizer words supported on its
complement. We now prescribe \(R\) and \(C\) and ask which full
stabilizers can have these two spaces.
Fix an isotropic subspace \(R\leq V_T\) of dimension \(r\), where
\(0\leq r\leq m\), and define
\[
 E=R_T^\perp/R,\qquad a=\dim E=2m-2r.
\]
Here \(R_T^\perp\) is the orthogonal complement of \(R\) within
\(V_T\).
An element of \(E\) is a coset \(v+R\): two local operators are
identified when their difference is a local stabilizer. This identification
retains their possible logical action, while the function below retains
their least physical weight.
The quotient carries a nondegenerate symplectic form. Fix a symplectic basis with matrix
\(J=\left[\begin{smallmatrix}0&I\\I&0\end{smallmatrix}\right]\).
For \(e\in E\), define
\[
 \cost(e)=\min_{v+R=e}\wt(v).
\]
\(c_R\) is the minimum coset weight in the original coordinate block.
An arbitrary
symplectic change of basis need not preserve it.

The dimension required of \(C\) follows from projection and rank--nullity.
For an actual extension, a vector on \(T\) is orthogonal to
\(\pi_T S^\perp\) if and only if it lies in \(S\cap V_T=R\).
\(\pi_T S^\perp=R_T^\perp\), of dimension \(2m-r\).
Its projection kernel is \(C\), so
\(\dim C=(n+k)-(2m-r)\).
Fix an additive code \(C\leq W\) of binary dimension
\(n+k-2m+r\), assumed to lie between zero and \(2f\). Put
\[
 P=C^\perp_W,\qquad b=\dim P=n-k-r.
\]
Orthogonal complements of \(C\) and \(P\) are taken inside \(W\).
Choose a basis \(v_1,\ldots,v_b\) of \(P\), set
\(\Gamma_{ij}=\langle v_i,v_j\rangle\), and identify \(W/C\)
with \(\F^b\) by the surjective syndrome map
\[
 \sigma(w)_j=\langle v_j,w\rangle,\qquad
 \delta(u)=\min\{\wt(w):\sigma(w)=u\}.
\]
Its kernel is \(C\). In particular \(\delta(0)=0\), regardless of
\(d(C)\). A syndrome \(u\) is a label for one coset of
\(C\), and \(\delta(u)\) is the smallest weight of any word with
that label.

For a subspace \(U\), its radical is
\(\rad(U)=\{u\in U:\langle u,v\rangle=0\text{ for every }v\in U\}\).
For a code, this same intersection with its dual is called its hull.
The nullspace of \(\Gamma\) represents \(\rad(P)=P\cap P^\perp
=C^\perp\cap C\). The binary nullity of \(\Gamma\) is
\(h=\dim(C\cap C^\perp)=\dim\Hull(C)\).

For reference, the four spaces used in the extension criterion are:
\begin{center}\small
\begin{tabularx}{\linewidth}{@{}lXl@{}}\toprule
Space&Meaning&Binary dimension\\\midrule
\(R\)&Stabilizer words supported on \(T\)&\(r\)\\
\(E\)&Local operators commuting with \(R\), modulo \(R\)&\(a=2m-2r\)\\
\(C\)&Centralizer shortened on \(T\)&\(n+k-2m+r\)\\
\(P=C^\perp\)&Projections of stabilizer words onto the complement&\(b=n-k-r\)\\\bottomrule
\end{tabularx}
\end{center}

\subsection{Reconstructing the full stabilizer}

The matrix \(M\) specifies which suffix coset accompanies each local
class. Its entries alone do not guarantee a code: the two coordinate
blocks must give commuting stabilizer words, and their combined weights
must satisfy the distance requirement. These are the two conditions in
the following theorem.

\begin{theorem}[Stabilizer extension criterion]\label{thm:extension}
Assume \(0<k<n\) and \(d(C)\geq d\), with \(R,C\) and their dimensions
as above.
There is a bijection between:
\begin{enumerate}[label=(\roman*)]
\item isotropic stabilizers \(S\) of dimension \(n-k\) satisfying
\[
 S\cap V_T=R,\quad S^\perp\cap W=C,\quad
 \wt(v)\geq d\ \text{for every }v\in S^\perp\setminus R;
\]
\item matrices \(M\in\F^{b\times a}\) satisfying
\begin{equation}\label{eq:exact}
 \Gamma=MJM^\tr,\qquad
 \cost(e)+\delta(Me)\geq d\quad(0\ne e\in E).
\end{equation}
\end{enumerate}
The stabilizer corresponding to \(M\) is \(S_M=D_M^\perp\), where
\begin{equation}\label{eq:graph-D}
 D_M=\{(v,w):v\in R_T^\perp,\ \sigma(w)=M(v+R)\}.
\end{equation}
The weight hypothesis in (i) is equivalent to distance at least \(d\)
together with localization of every stabilizer word of weight below \(d\)
inside \(R\).
\end{theorem}

\begin{proof}
\emph{Recovering the attachment map.}
Let \(S\) satisfy (i), and put \(D=S^\perp\). The annihilator of
\(\pi_TD\) inside \(V_T\) is \(S\cap V_T=R\); hence
\(\pi_TD=R_T^\perp\). The kernel of \(\pi_T|_D\) is \(C\).
For each \(v\in R_T^\perp\), the corresponding words on the
complement of \(T\) form a nonempty coset of \(C\). Addition in \(D\) makes this coset assignment linear. Since \(R\subseteq S\subseteq D\), it depends only
on \(e=v+R\), defining \(M:E\to W/C\).

For \(e\ne0\), its prefix representatives and the members of its suffix
coset can be chosen independently, by adding \(R\) and \(C\).
Their minimum total weight is
\(\cost(e)+\delta(Me)\). All lie outside \(R\), proving the inequality.
The suffix-only fiber gives \(d(C)\geq d\).

\emph{Recovering the commutation condition.}
The duality between projection and shortening also gives \(\pi_WS=P\). Indeed,
a word on \(W\) annihilates \(\pi_WS\) if and only if its extension
by zeros belongs to \(S^\perp\cap W=C\), so
\((\pi_WS)^\perp_W=C\). The kernel of
\(\pi_W|_S\) is \(R\), so \(S/R\) is the graph of a linear map
\(\tau:P\to E\). Write \(A\) for its \(a\)-by-\(b\) coordinate matrix.
For \(x=\sum c_jv_j\), this graph has prefix class \(Ac\).
Commutation with a centralizer word of prefix class \(e\) gives
\(e^\tr JAc+c^\tr Me=0\) for every \(e,c\). Thus
\(M=A^\tr J\), equivalently \(A=JM^\tr\). Two graph words
commute if and only if their local and complementary symplectic products
cancel. Isotropy of the graph gives
\(\Gamma=A^\tr JA=MJM^\tr\). \(\tau\) is the symplectic
adjoint of \(M\).
Its kernel has a direct interpretation: \(\tau(x)=0\) means that
a stabilizer projecting to \(x\) has local part in \(R\), which can
be removed by adding that local stabilizer. \(\ker\tau=S\cap W\).

\emph{Constructing a stabilizer from the map.}
Conversely, let \(M\) satisfy (ii) and define \(D_M\) by
\eqref{eq:graph-D}. Its projection is \(R_T^\perp\), its suffix-only
kernel is \(C\), and it contains \(R\). Consequently
\(\dim D_M=2m-r+\dim C=n+k\), so \(\dim S_M=n-k\),
and annihilator duality gives \(S_M\cap V_T=R\).
For a projected stabilizer word \(x=\sum c_jv_j\), commutation with
\(D_M\) forces its prefix class to be \(JM^\tr c\). The symplectic product of
two such words is
\(c^\tr(MJM^\tr+\Gamma)c'=0\). The local words of \(R\) commute
with them because their prefixes lie in \(R_T^\perp\). \(S_M\)
is isotropic.

Every fiber of nonzero prefix class has weight at least \(d\) by (ii).
In the zero-class fiber, the prefix lies in \(R\) and the suffix lies in
\(C\); a word outside \(R\) has a nonzero suffix, of weight at least
\(d\). This proves (i). The fibers determine \(D\) and therefore
\(S=D^\perp\) uniquely, establishing the bijection.
\end{proof}

\subsection{A relaxation using only coset weights}

The distance inequalities do not generally imply commutation. For example,
take \(n=4,k=1,d=2\), \(R=\langle ZZ\rangle\), and
\(C=\langle XX,ZZ\rangle\) on the other two sites. In the prefix basis
\(Z_0,X_0X_1\), the nonzero coset minimum weights are \(1,2,2\). Choose independent
nonzero syndrome columns. All weight inequalities hold, since the three
nonzero cosets of \(C\) have distance one. But \(C\) is self-dual,
so \(\Gamma=0\), whereas \(MJM^\tr\ne0\).
Excluding all weight-admissible maps still excludes genuine extensions.
Finding one such map would not construct a code.

\subsection{Extensions of degenerate codes}

For an arbitrary degenerate code, \(C\) may contain short stabilizer words.
The same graph construction still applies, with the distance taken relative
to the reconstructed stabilizer.

\begin{theorem}[Relative extension criterion]\label{thm:relative}
Keep \(R,C,E,P\) and their dimensions, but impose no lower bound on
\(d(C)\). Let \(M\) satisfy \(\Gamma=MJM^\tr\), and define
\(\tau_M(\sum c_jv_j)=JM^\tr c\), and
\[
 \delta_M^{\rm log}(e)=\min\{\wt(w):\sigma(w)=Me,\
             w\notin\tau_M^{-1}(e)\},
\]
where \(\tau_M^{-1}(e)\) is a subset of \(P\), and the minimum of
the empty set is \(+\infty\).
Then \(S_M\) has distance at least \(d\) if and only if
\[
 \cost(e)+\delta_M^{\rm log}(e)\geq d\quad\text{for every }e\in E,
\]
including \(e=0\).
\end{theorem}
\begin{proof}
The graph and dimension arguments above do not require a distance assumption.
A word \((v,w)\in D_M\), with \(v+R=e\), is in \(S_M\) if and only if
\(w\in P\) and \(\tau_M(w)=e\). Membership is independent of the
representative of \(e\), because \(R\subseteq S_M\). The displayed
minimum ranges over precisely the suffixes of logical words in
that fiber. Prefix and suffix representatives vary independently. Minimizing
their sum gives the assertion, fiber by fiber.
\end{proof}

At \(e=0\), the suffix minimum is the distance of \(C\) relative to
\(N=\ker\tau_M\): it is \(\min\{\wt(w):w\in C\setminus N\}\).
The coset minimum depends on \(M\), so the theorem does not supply a fixed
precomputed table of coset minimum weights for arbitrary degeneracy. The searches below
use Theorem~\ref{thm:extension} only after proving purity or localization.
